\begin{proposition}[5.1.2]
\label{IV.2.5.1.2-fr}
Sean $X$ un preesquema, $Y$ una parte cerrada irreducible de $X$, e $y$ el
punto genérico de $Y$. Se tiene
\begin{equation}
\operatorname{codim}(Y,X)=\dim(\mathcal O_{X,y})
\tag{5.1.2.1}\label{IV.2.5.1.2.1-fr}
\end{equation}
\end{proposition}

\begin{proof}
En efecto (\textbf I, 2.4.2), las partes cerradas irreducibles de $X$ que contienen a
$y$ están canónicamente en correspondencia biunívoca con las partes cerradas
irreducibles de $\operatorname{Spec}(\mathcal O_{X,y})$ y, por tanto, están en
correspondencia biunívoca con los ideales primos de $\mathcal O_{X,y}$, y
(5.1.2.1) resulta de las definiciones.

En particular, se obtiene así de nuevo el hecho de que los puntos genéricos de
las componentes irreducibles de $X$ son los \emph{únicos} puntos $x\in X$ tales
que $\dim(\mathcal O_x)=0$ (\textbf I, 1.1.14).
\end{proof}

\oldpage[IV-2]{87}

\begin{corollary}[5.1.3]
\label{IV.2.5.1.3-fr}
Para toda parte cerrada $Y$ de un preesquema $X$, se tiene
\begin{equation}
\operatorname{codim}(Y,X)=\inf_{y\in Y}\dim(\mathcal O_{X,y}).
\tag{5.1.3.1}\label{IV.2.5.1.3.1-fr}
\end{equation}
Además, si $X$ es localmente noetheriano, se tiene, para todo $x\in Y$,
\begin{equation}
\operatorname{codim}_x(Y,X)=
\inf_{\substack{y\in Y,\,x\in\overline{\{y\}}}}
\dim(\mathcal O_{X,y}).
\tag{5.1.3.2}\label{IV.2.5.1.3.2-fr}
\end{equation}
\end{corollary}

\begin{proof}
La relación (5.1.3.2) resulta, en efecto, de (5.1.3.1) y de $(\mathbf{0},14.2.6)$.

Este corolario permite \emph{definir}, para una parte \emph{arbitraria} $Y$ de
un preesquema $X$, la \emph{codimensión} $\operatorname{codim}(Y,X)$ de $Y$ en
$X$ como igual al segundo miembro de (5.1.3.1).
\end{proof}

\begin{proposition}[5.1.4]
\label{IV.2.5.1.4-fr}
Para todo preesquema $X$, se tiene
\begin{equation}
\dim(X)=\sup_{x\in X}(\dim(\mathcal O_x)).
\tag{5.1.4.1}\label{IV.2.5.1.4.1-fr}
\end{equation}
Si toda parte cerrada irreducible de $X$ contiene un punto cerrado, se tiene
también
\begin{equation}
\dim(X)=\sup_{x\in F}(\dim(\mathcal O_x))
\tag{5.1.4.2}\label{IV.2.5.1.4.2-fr}
\end{equation}
donde $F$ es el conjunto de los puntos cerrados de $X$.
\end{proposition}

\begin{proof}
Basta observar (en virtud de (\textbf I, 2.4.2)) que las cadenas de partes cerradas
irreducibles de $X$ corresponden biunívocamente a las cadenas de partes
cerradas irreducibles de un esquema local de $X$; cuando toda parte cerrada
irreducible de $X$ contiene un punto cerrado, es posible evidentemente limitarse
a los esquemas locales en los puntos cerrados de $X$.

Se observará que toda parte cerrada irreducible de $X$ contiene un punto cerrado
cuando $X$ es \emph{cuasicompacto} $(\mathbf{0}_{\mathrm I},2.1.3)$; veremos un poco más
adelante (5.1.11) que lo mismo ocurre cuando $X$ es \emph{localmente
noetheriano}.
\end{proof}

\begin{corollary}[5.1.5]
\label{IV.2.5.1.5-fr}
Para que un preesquema $X$ sea catenario, es necesario y suficiente que, para
todo $x\in X$, el anillo local $\mathcal O_x$ sea catenario. Si, además, toda
parte cerrada irreducible de $X$ contiene un punto cerrado, basta, para que $X$
sea catenario, con que $\mathcal O_x$ sea catenario para todo punto cerrado $x$
de $X$.
\end{corollary}

\begin{proof}
El razonamiento es el mismo que en (5.1.4), puesto que se trata de comparar
cadenas de partes cerradas irreducibles que tienen los mismos extremos.
\end{proof}

